\newgeometry{top=2.5cm, bottom=2.5cm, left=2cm, right=2cm}
\setlength{\parindent}{0pt}
\setlength{\parskip}{6pt}

% --- TITOLO ---
\begin{center}
    {\Huge \textbf{\color{cherryRed} Vettori Applicati nel Piano}} \\
    \vspace{0.3cm}
    {\Large \textit{Struttura dello Spazio Vettoriale $V_O^2$}}
    \vspace{0.5cm}
    \hrule height 1pt
\end{center}

\section{Costruzione Geometrica di $V_O^2$}

Dalla proiezione in aula (slide 3), costruiamo lo spazio vettoriale basato sulla geometria euclidea.

\begin{definition}{Vettore Applicato in $O$}{}
Consideriamo il piano euclideo $\mathbb{R}^2$ e fissiamo un punto $O \in \mathbb{R}^2$ detto \textbf{Origine} (o punto fissato).

Un \textbf{Vettore Applicato in $O$} è definito come un segmento orientato con:
\begin{itemize}
    \item \textbf{Primo estremo:} il punto $O$.
    \item \textbf{Secondo estremo:} un punto $A \in \mathbb{R}^2$.
\end{itemize}
Si denota con $\overrightarrow{OA}$.
\end{definition}

\subsection{L'Insieme $V_O^2$}
Definiamo l'insieme di tutti i vettori applicati nell'origine:
\[
V_O^2 \doteq \{ \text{vettori applicati in } O \}
\]
L'\textbf{Origine} di $V_O^2$ corrisponde al vettore nullo, ovvero il punto $O$ stesso (segmento degenere).

% =========================================================================
\section{Operazioni in $V_O^2$}

Per dotare l'insieme $V_O^2$ di una struttura di spazio vettoriale, definiamo le due operazioni fondamentali.

\subsection{1. Somma di Vettori}
L'operazione somma è definita come:
\[ + : V_O^2 \times V_O^2 \longrightarrow V_O^2 \]
Dati due vettori $\overrightarrow{OA}$ e $\overrightarrow{OB}$, la somma è il vettore $\overrightarrow{OC}$ tale che $C$ è il quarto vertice del parallelogramma costruito su $OA$ e $OB$.

\[ \overrightarrow{OA} + \overrightarrow{OB} = \overrightarrow{OC} \]

\begin{center}
\begin{tikzpicture}[>=Stealth, scale=1.5]
    % Coordinate
    \coordinate (O) at (0,0);
    \coordinate (A) at (2, 0.5);
    \coordinate (B) at (1, 1.5);
    \coordinate (C) at (3, 2); % A + B

    % Vettori
    \draw[->, thick, color=maroon] (O) -- (A) node[midway, below right] {$\overrightarrow{OA}$};
    \draw[->, thick, color=maroon] (O) -- (B) node[midway, left] {$\overrightarrow{OB}$};
    \draw[->, ultra thick, color=cherryRed] (O) -- (C) node[midway, above left] {$\overrightarrow{OC} = \text{Somma}$};

    % Linee tratteggiate per parallelogramma
    \draw[dashed, gray] (A) -- (C);
    \draw[dashed, gray] (B) -- (C);

    % Etichette Punti
    \node[below left] at (O) {$O$};
    \node[right] at (A) {$A$};
    \node[above] at (B) {$B$};
    \node[above right] at (C) {$C$};
\end{tikzpicture}
\end{center}

\subsection{2. Prodotto per Scalare}
L'operazione prodotto per uno scalare $\lambda \in \mathbb{R}$ è definita come:
\[ \cdot : \mathbb{R} \times V_O^2 \longrightarrow V_O^2 \]
Dato uno scalare $\lambda$ e un vettore $\overrightarrow{OA}$, il risultato è un vettore $\overrightarrow{OC}$ tale che:
\begin{itemize}
    \item $C$ appartiene alla retta passante per $O$ e per $A$.
    \item La lunghezza (modulo) soddisfa:
    \[ \frac{\text{lunghezza } OC}{\text{lunghezza } OA} = |\lambda| \]
\end{itemize}
Se $\lambda > 0$, il verso è lo stesso. Se $\lambda < 0$, il verso è opposto.

\textbf{Esempio dalla lavagna ($\lambda = 2$):}

\begin{center}
\begin{tikzpicture}[>=Stealth, scale=1.5]
    % Coordinate
    \coordinate (O) at (0,0);
    \coordinate (A) at (1.5, 0.5);
    \coordinate (C) at (3, 1); % 2 * A

    % Retta tratteggiata
    \draw[dashed, gray] (-0.5, -0.16) -- (3.5, 1.16);

    % Vettori
    \draw[->, thick, color=maroon] (O) -- (A) node[midway, above left] {$\overrightarrow{OA}$};
    \draw[->, ultra thick, color=cherryRed, transform canvas={yshift=-2pt}] (O) -- (C) node[midway, below right] {$\lambda \cdot \overrightarrow{OA}$};

    % Etichette
    \node[below left] at (O) {$O$};
    \node[above] at (A) {$A$};
    \node[above] at (C) {$C$};
    \node[right] at (3.5, 1) {$\lambda=2$};
\end{tikzpicture}
\end{center}

\begin{tcolorbox}[colback=white, colframe=cherryRed, title=Affermazione (AFF)]
L'insieme $V_O^2$, dotato delle operazioni di somma (regola del parallelogramma) e prodotto per scalare (dilatazione/contrazione), è uno \textbf{Spazio Vettoriale} su $\mathbb{R}$.
\end{tcolorbox}

\begin{observation}{Scopo della Lezione (Blackboard)}{}
Alla base della lavagna si legge:
\textit{"Scopo: Dotare $V_O^2$ di una struttura di spazio vettoriale."} 
Questo significa verificare che le operazioni geometriche appena definite soddisfino gli 8 assiomi visti precedentemente (associatività, elemento neutro, ecc.).
\end{observation}

\restoregeometry
